\chapter*{Conclusion}
\addcontentsline{toc}{chapter}{Conclusion}
\markboth{Conclusion}{Conclusion}

The three essays make visible a problem that neither capacity measurement nor a monetary account of crisis can settle alone: the reproduction of production depends on relations that a fixed technical benchmark or an autonomous monetary impulse leaves unexplained. Identifying productive capacity requires an account of how it is formed. Explaining its formation requires attention to accumulation, distribution, technique, and the organization of production. In the periphery, sustaining production also requires access to foreign exchange and international settlement. These connections give the dissertation its unity while preserving the difference between a long-run structural relation and a historically specific crisis.

\section*{From the Capacity Benchmark to the Conditions of Reproduction}

Chapter~1 establishes the empirical difficulty of separating the productive-capacity denominator from the relations governing its formation. In the full-sample US replication, the bivariate output--capital system does not sustain the required long-run properties. The surviving relation incorporates the rate of exploitation and historical structure. The measurement problem is therefore substantive: changing the treatment of distribution changes the basis on which capacity can be identified. A fitted output path alone cannot establish an autonomous technical ceiling.

Chapter~2 develops the structural account left open by that result. Accumulated capital becomes usable capacity through choices of technique made by decentralized and internally divided firms. Distribution affects those choices, and the composition of the capital stock conditions the resulting transformation. The estimated elasticities below unity mean that capacity formation does not keep pace proportionally with accumulation within the chapter's specifications. They do not mean that productive capacity must fall in absolute terms or that a crisis follows at a determinate date. The distinction between capacity formation and realized utilization also matters: a changing transformation elasticity does not itself establish how much of the resulting capacity will be used.

Reading the first two essays together thus changes the status of the capacity benchmark. It becomes a relation to be explained and estimated, rather than an engineering constant that can be imported into an analysis of demand or profitability. Functional distribution, labor-process conflict, and capital composition enter this account at different levels. Their connection supplies an interpretation of the macroeconomic relation without making an aggregate coefficient a direct estimate of every organizational process inside the firm.

\section*{Peripheral Accumulation, International Money, and Agency}

The Chilean extension introduces a material limit to the conversion of an investment incentive into productive capacity. When machinery must be imported, access to foreign exchange conditions whether firms can act on a distributive incentive to mechanize. Chapter~2's sharply attenuated mechanization response in the constrained state places the balance-of-payments constraint inside the transformation process. The center--periphery comparison therefore concerns the conditions under which accumulation can proceed. Different national specifications and historical samples limit direct numerical comparison, but they do not erase the analytical importance of import dependence.

Chapter~3 shows why the explanation of peripheral crisis cannot stop at that long-run relation. The Unidad Popular confronted the reproduction of enterprise liquidity and international payments as well as productive capacity. The open-economy balance-sheet framework connects domestic monetary liabilities to reserve assets denominated in internationally acceptable means of settlement. Financial subordination and currency hierarchy describe the broader international relations; the solvency measure describes a specific balance-sheet condition within them. Domestic monetary authority does not, by itself, provide command over world money.

The historical analysis distinguishes global monetary instability following the breakdown of Bretton Woods, commodity and import-price movements, targeted bilateral and commercial financial pressure, and domestic distributive and political conflict. Their interaction narrowed the BCCh's policy space without determining one inevitable response. The bank's choice to accommodate enterprise liquidity needs remains an institutional action under constraint. The dynamic estimates provide evidence consistent with that interpretation: in the adverse solvency-growth state, reverse accommodation is more persistent and cumulatively larger than forward monetary transmission over the reported horizon. Forward transmission remains present. Endogenous accommodation therefore qualifies an account centered on autonomous money creation without implying that monetary expansion has no inflationary consequences.

Together, the essays distinguish two roles of the external constraint. It can restrict the acquisition of the means of production, and it can condition the monetary and financial processes required to sustain ongoing production. The long-run mechanization threshold and the monthly solvency-growth threshold address those different questions. Neither should be read as the other's statistical validation. Their connection is an argument about the conditions of peripheral reproduction, not one universal crisis law spanning all three essays.

\section*{Limits and Further Inquiry}

The dissertation's limits follow from this division of evidence. Chapter~1 uses an aggregate, single-sector framework and a time-invariant transformation parameter. Chapter~2 extends the structural account but does not supply a complete theory of crisis or profitability, and its national specifications are not fully harmonized. Chapter~3 establishes directional precedence and state-dependent conditional responses rather than experimental counterfactual identification. Unobserved common shocks, historical breaks, and the distance between macroeconomic aggregates and institutional decisions remain consequential for interpretation.

Two extensions would deepen the connections developed here. More comparable capital-composition series for the United States and Chile would allow the transformation relation to be estimated on a more common basis. Historical evidence tracing credit allocation between the BCCh, commercial banks, and enterprises would connect the aggregate accommodation patterns more closely to institutional decisions. Comparative research could then examine whether the Chilean solvency dynamics recur elsewhere, without assuming that the estimated threshold transfers unchanged.

Capitalist reproduction cannot be understood by assuming that accumulated capital, productive capacity, and macroeconomic stability expand in step. Connecting capacity measurement, structural explanation, and historical monetary analysis, while retaining the limits of their evidence, is the dissertation's contribution to the political economy of accumulation and crisis in center and periphery.
